\begin{recipe}{Kip madras}
\kicker[Hoofdgerecht,Vlees,Indiaas,Doordeweeks]{Hoofdgerecht · vlees · Indiaas · doordeweeks}
\index[register]{Kip madras}
\index[register]{Kip}
\index[register]{Bloemkool}
\index[register]{Kokosmelk}

\meta{35 minuten}

\lettrine{E}{en} avondje lui doordeweeks. Dit recept is gebouwd rond het
verspakket curry madras van de Lidl, met een extra zakje voorgesneden
bloemkool erbij voor wat meer volume. Gebakken uitjes geven de curry wat
bite en de sambal staat los op tafel voor wie van pittig houdt.

\blockrule{Ingrediënten}
\begin{ingredients}
  \ing{1}{kipfilet, in stukken}\index[register]{Kip}
  \ing{170 g}{langkorrelige rijst}
  \ing{400 g}{bloemkool, in roosjes (extra)}\index[register]{Bloemkool}
  \ing{1}{verspakket curry madras (Lidl)}
  \ingb{olijfolie}
  \ingb{gebakken uitjes}
  \ingb{sambal, voor de liefhebber}
\end{ingredients}

\blockrule{Bereiding}
\tip{Dit werkt ook met een ander verspakket curry. Zolang je zelf kip en
rijst toevoegt kun je de verpakking gewoon volgen.}
\begin{steps}
  \item Kook de rijst volgens de aanwijzingen op de verpakking.
  \item Maak de groenten uit het verspakket schoon en snijd ze zoals op de
        verpakking staat aangegeven.
  \item Bak de kipfilet aan in een scheut olijfolie en volg verder de
        bereidingsstappen van het verspakket.
  \item Voeg de extra zak bloemkool samen met de rest van de groenten toe,
        zodat er meer bij is dan het pakket alleen biedt.
  \item Laat de curry volgens de verpakking gaar stoven.
  \item Serveer de kip madras met de rijst, gebakken uitjes en sambal
        erbij.
\end{steps}
\end{recipe}
